% Metodología
\section{Metodología}

\begin{frame}
\frametitle{Datos y fuentes — 5 módulos integrados}
\begin{columns}[T]
\column{0.52\textwidth}
\textbf{Fuentes}
\begin{itemize}
\item Base limpia de consumo (cap. K)
\item Módulo \texttt{personas} (sexo, edad)
\item Módulo \texttt{d2} (educación, hijos)
\item Módulo \texttt{g} (red social, actitud)
\item Módulo \texttt{d} (laboral, salud, percepción)
\end{itemize}

\column{0.44\textwidth}
\textbf{Cobertura}
\begin{itemize}
\item \(n=3\,982\) individuos
\item 16 features, 0\% nulos (salvo \texttt{regimen\_salud}: 11.2\%)
\item \texttt{consumo\_12m}: 1\,719 positivos
\item Precio: submuestra exploratoria ($n{\approx}755$)
\end{itemize}
\end{columns}
\vspace{0.3cm}
\textit{Decisión}: ML sobre la muestra amplia sin precio.
\end{frame}

\begin{frame}
\frametitle{Pipeline reproducible}
\begin{enumerate}
\item \texttt{python -m cannabis\_tax.cli pipeline}: limpieza, validación y construcción de la base modeladora.
\item Notebook 01: EDA + benchmark MCO/Probit.
\item Notebook 03: entrenamiento RF, XGBoost, GBM.
\item Notebook 05: CV-10, búsqueda de hiperparámetros, selección del modelo final.
\item Notebook 06: SHAP + robustez + ética.
\end{enumerate}
\vspace{0.3cm}
\textit{Tests}: \texttt{make test} corre la suite unitaria.
\end{frame}

\begin{frame}
\frametitle{Benchmark econométrico}
\textbf{MCO (LPM)}:
\[
y_i = \beta_0 + \beta_2 edad_i + \beta_3 edad_i^2 + \beta_4 mujer_i + \beta_5 educacion_i + u_i
\]
con errores HC1.

\vspace{0.3cm}
\textbf{Probit}:
\[
\Pr(y_i=1\mid X_i) = \Phi(\alpha_0 + \alpha_2 edad_i + \alpha_3 edad_i^2 + \alpha_4 mujer_i + \alpha_5 educacion_i)
\]
con efectos marginales promedio.

\vspace{0.3cm}
\textit{Roles}: referencia interpretable y punto de comparación.
\end{frame}

\begin{frame}
\frametitle{Feature set ampliado — 16 variables}
\centering
\footnotesize
\begin{tabular}{ll}
\toprule
\textbf{Grupo} & \textbf{Variables} \\
\midrule
Demográficas base & edad, edad², sexo, educación \\
Estructura familiar & n\_hijos \\
Red social (G\_01, G\_02) & familiares\_consumen, amigos\_consumen \\
Actitud (G\_04, G\_09) & probaria\_sustancias, cannabis\_medicinal \\
Situación socioeconómica (D\_02, D\_07) & actividad\_laboral, regimen\_salud \\
Salud mental PHQ-2 (D\_09, D\_10) & sintomas\_depresivos, anhedonia \\
Percepción de riesgo (D\_11\_F, D\_11\_G) & riesgo\_marihuana\_ocasional/frecuente \\
\bottomrule
\end{tabular}
\vspace{0.3cm}
\\
\small Todos con 0\% de nulos salvo \texttt{regimen\_salud} (11.2\%). Categóricas con \textit{one-hot encoding}.
\end{frame}

\begin{frame}
\frametitle{Modelos de Machine Learning}
\begin{itemize}
    \item \textbf{Random Forest}: \textit{bagging}, 300 árboles, \texttt{max\_depth=7}.
    \item \textbf{XGBoost}: \textit{boosting} secuencial, ponderado por prevalencia.
    \item \textbf{Gradient Boosting (sklearn)}: \textit{boosting} conservador, \texttt{max\_depth=3}.
\end{itemize}
\vspace{0.3cm}
\textbf{Reglas de comparación}:
\begin{itemize}
    \item \textit{Feature set} ampliado de 16 variables (ver tabla anterior).
    \item Mismo \texttt{train/test split} estratificado (\texttt{random\_state=42}).
    \item Métricas: Accuracy, AUC-ROC, F1, CV AUC (5 y 10 folds).
\end{itemize}
\end{frame}

\begin{frame}
\frametitle{Validación y selección}
\begin{enumerate}
    \item \textbf{CV-10}: estabilidad y solapamiento entre modelos.
    \item \textbf{Curvas de aprendizaje}: control de \textit{overfitting}.
    \item \textbf{\texttt{RandomizedSearchCV}}: 60 combinaciones × 5 folds sobre el GBM.
    \item \textbf{Criterios}: AUC-ROC, F1, brecha train-validation, plausibilidad económica.
\end{enumerate}
\vspace{0.3cm}
\textit{No se evalúa sólo precisión}: la argumentación pesa tanto como el número.
\end{frame}
